\documentclass[10pt,a4paper]{amsart}
\usepackage{amsmath,amssymb,amsthm,graphicx}
\usepackage{verbatim}
\newenvironment{reference}{\comment}{\endcomment}
%\newenvironment{reference}{}{}
\newenvironment{slogan}{\comment}{\endcomment}
\newenvironment{history}{\comment}{\endcomment}

% For commutative diagrams we use Xy-pic
\usepackage[all]{xy}

% We use 2cell for 2-commutative diagrams.
\xyoption{2cell}
% Theorem environments.
%
\theoremstyle{plain}
\newtheorem{theorem}[subsection]{Theorem}
\newtheorem{proposition}[subsection]{Proposition}
\newtheorem{lemma}[subsection]{Lemma}

\theoremstyle{definition}
\newtheorem{definition}[subsection]{Definition}
\newtheorem{example}[subsection]{Example}
\newtheorem{exercise}[subsection]{Exercise}
\newtheorem{situation}[subsection]{Situation}

\theoremstyle{remark}
\newtheorem{remark}[subsection]{Remark}
\newtheorem{remarks}[subsection]{Remarks}

\numberwithin{equation}{subsection}

% Macros
%
\def\lim{\mathop{\mathrm{lim}}\nolimits}
\def\colim{\mathop{\mathrm{colim}}\nolimits}
\def\Spec{\mathop{\mathrm{Spec}}}
\def\Hom{\mathop{\mathrm{Hom}}\nolimits}
\def\Ext{\mathop{\mathrm{Ext}}\nolimits}
\def\SheafHom{\mathop{\mathcal{H}\!\mathit{om}}\nolimits}
\def\SheafExt{\mathop{\mathcal{E}\!\mathit{xt}}\nolimits}
\def\Sch{\mathit{Sch}}
\def\Mor{\mathop{\mathrm{Mor}}\nolimits}
\def\Ob{\mathop{\mathrm{Ob}}\nolimits}
\def\Sh{\mathop{\mathit{Sh}}\nolimits}
\def\NL{\mathop{N\!L}\nolimits}
\def\CH{\mathop{\mathrm{CH}}\nolimits}
\def\proetale{{pro\text{-}\acute{e}tale}}
\def\etale{{\acute{e}tale}}
\def\QCoh{\mathit{QCoh}}
\def\Ker{\mathop{\mathrm{Ker}}}
\def\Im{\mathop{\mathrm{Im}}}
\def\Coker{\mathop{\mathrm{Coker}}}
\def\Coim{\mathop{\mathrm{Coim}}}

% Boxtimes
%
\DeclareMathSymbol{\boxtimes}{\mathbin}{AMSa}{"02}

%
% Macros for moduli stacks/spaces
%
\def\QCohstack{\mathcal{QC}\!\mathit{oh}}
\def\Cohstack{\mathcal{C}\!\mathit{oh}}
\def\Spacesstack{\mathcal{S}\!\mathit{paces}}
\def\Quotfunctor{\mathrm{Quot}}
\def\Hilbfunctor{\mathrm{Hilb}}
\def\Curvesstack{\mathcal{C}\!\mathit{urves}}
\def\Polarizedstack{\mathcal{P}\!\mathit{olarized}}
\def\Complexesstack{\mathcal{C}\!\mathit{omplexes}}
% \Pic is the operator that assigns to X its picard group, usage \Pic(X)
% \Picardstack_{X/B} denotes the Picard stack of X over B
% \Picardfunctor_{X/B} denotes the Picard functor of X over B
\def\Pic{\mathop{\mathrm{Pic}}\nolimits}
\def\Picardstack{\mathcal{P}\!\mathit{ic}}
\def\Picardfunctor{\mathrm{Pic}}
\def\Deformationcategory{\mathcal{D}\!\mathit{ef}}
\usepackage{cleveref}
\crefname{originalequation}{original equation}{original equations}
\Crefname{originalequation}{Original equation}{Original equations}
\makeatletter
\newlabel{equation-FM-QCoh}{{0FZC}{source}{}{sourcepointer-1}{}}
\newlabel{equation-FM-QCoh@cref}{{[originalequation][0][]0FZC}{[1][0][]source}{}{}{}}
\makeatother
\crefname{originalreference}{original reference}{original references}
\Crefname{originalreference}{Original reference}{Original references}
\makeatletter
\newlabel{morphisms-lemma-affine-permanence}{{01SG}{source}{}{sourcepointer-2}{}}
\newlabel{morphisms-lemma-affine-permanence@cref}{{[originalreference][0][]01SG}{[1][0][]source}{}{}{}}
\makeatother
\crefname{originalreference}{original reference}{original references}
\Crefname{originalreference}{Original reference}{Original references}
\makeatletter
\newlabel{coherent-lemma-relative-affine-vanishing}{{01XC}{source}{}{sourcepointer-3}{}}
\newlabel{coherent-lemma-relative-affine-vanishing@cref}{{[originalreference][0][]01XC}{[1][0][]source}{}{}{}}
\makeatother
\crefname{originalreference}{original reference}{original references}
\Crefname{originalreference}{Original reference}{Original references}
\makeatletter
\newlabel{coherent-lemma-colimit-cohomology}{{07TB}{source}{}{sourcepointer-4}{}}
\newlabel{coherent-lemma-colimit-cohomology@cref}{{[originalreference][0][]07TB}{[1][0][]source}{}{}{}}
\makeatother
\crefname{originallemma}{original lemma}{original lemmas}
\Crefname{originallemma}{Original lemma}{Original lemmas}
\makeatletter
\newlabel{lemma-functor-quasi-coherent-from-affine}{{0FZD}{source}{}{sourcepointer-5}{}}
\newlabel{lemma-functor-quasi-coherent-from-affine@cref}{{[originallemma][0][]0FZD}{[1][0][]source}{}{}{}}
\makeatother
\crefname{originallemma}{original lemma}{original lemmas}
\Crefname{originallemma}{Original lemma}{Original lemmas}
\makeatletter
\newlabel{lemma-functor-quasi-coherent-from-affine-compose}{{0FZF}{source}{}{sourcepointer-6}{}}
\newlabel{lemma-functor-quasi-coherent-from-affine-compose@cref}{{[originallemma][0][]0FZF}{[1][0][]source}{}{}{}}
\makeatother
\crefname{originalreference}{original reference}{original references}
\Crefname{originalreference}{Original reference}{Original references}
\makeatletter
\newlabel{coherent-lemma-separated-case-relative-cech}{{01XL}{source}{}{sourcepointer-7}{}}
\newlabel{coherent-lemma-separated-case-relative-cech@cref}{{[originalreference][0][]01XL}{[1][0][]source}{}{}{}}
\makeatother
\title{Constructing the kernel and a comparison transformation over quasi-compact schemes with affine diagonal}
\author{The Stacks Project Authors}
\date{stacks-project:functors.tex:local-snapshot-20261012; GNU FDL 1.2 or later}
\begin{document}
\maketitle
\noindent\textit{Editorial note.} All delivered mathematical text is the author's own, copied verbatim from the pinned Stacks Project chapter file functors.tex; no statement or proof step was rewritten, reordered or removed.\par\medskip
\noindent\textit{Editorial note.} Span recomputation by label pairing (LANE\_BRIEF 10.2): in this chapter the author's proof environment follows the statement environment as a standalone block; the spec covers the complete environment pair (statement span, proof span) as recorded in the pool candidate.\par\medskip
\noindent\textit{Editorial note.} Occupation scan: combined claim\_key + statement\_anchor + generic-line scan over all registration files (active and withdrawn) found 0 hits for functors.tex (fu-occupied.json, 2026-10-14); consistent with the duty manager's six-check dispatch.\par\medskip
\ifdefined\hypertarget\hypertarget{sourcepointer-1}{}\fi
\noindent\textit{Source reference.} Original equation 0FZC: sources/stacks-chow/tags.txt, source lines 19947--19947. Cross-chapter standard prerequisite cited by the author inside the counted span; it is used as a pointer and its own proof is not part of this item.\par\medskip
\ifdefined\hypertarget\hypertarget{sourcepointer-2}{}\fi
\noindent\textit{Source reference.} Original reference 01SG: sources/stacks-chow/tags.txt, source lines 2337--2337. Cross-chapter standard prerequisite cited by the author inside the counted span; it is used as a pointer and its own proof is not part of this item.\par\medskip
\ifdefined\hypertarget\hypertarget{sourcepointer-3}{}\fi
\noindent\textit{Source reference.} Original reference 01XC: sources/stacks-chow/tags.txt, source lines 2513--2513. Cross-chapter standard prerequisite cited by the author inside the counted span; it is used as a pointer and its own proof is not part of this item.\par\medskip
\ifdefined\hypertarget\hypertarget{sourcepointer-4}{}\fi
\noindent\textit{Source reference.} Original reference 07TB: sources/stacks-chow/tags.txt, source lines 9931--9931. Cross-chapter standard prerequisite cited by the author inside the counted span; it is used as a pointer and its own proof is not part of this item.\par\medskip
\ifdefined\hypertarget\hypertarget{sourcepointer-5}{}\fi
\noindent\textit{Source reference.} Original lemma 0FZD: sources/stacks-chow/tags.txt, source lines 19948--19948. Cross-chapter standard prerequisite cited by the author inside the counted span; it is used as a pointer and its own proof is not part of this item.\par\medskip
\ifdefined\hypertarget\hypertarget{sourcepointer-6}{}\fi
\noindent\textit{Source reference.} Original lemma 0FZF: sources/stacks-chow/tags.txt, source lines 19950--19950. Cross-chapter standard prerequisite cited by the author inside the counted span; it is used as a pointer and its own proof is not part of this item.\par\medskip
\ifdefined\hypertarget\hypertarget{sourcepointer-7}{}\fi
\noindent\textit{Source reference.} Original reference 01XL: sources/stacks-chow/tags.txt, source lines 2522--2522. Cross-chapter standard prerequisite cited by the author inside the counted span; it is used as a pointer and its own proof is not part of this item.\par\medskip
\begin{lemma}
\label{lemma-functor-quasi-coherent-from-affine-diagonal-pre}
Let $R$ be a ring. Let $X$ and $Y$ be schemes over $R$. Assume $X$
is quasi-compact with affine diagonal. Let
$F : \QCoh(\mathcal{O}_X) \to \QCoh(\mathcal{O}_Y)$
be an $R$-linear, right exact functor which commutes
with arbitrary direct sums. Then we can construct
\begin{enumerate}
\item a quasi-coherent module $\mathcal{K}$ on $X \times_R Y$, and
\item a natural transformation $t : F \to F_\mathcal{K}$
where $F_\mathcal{K}$ denotes the functor (\ref{equation-FM-QCoh})
\end{enumerate}
such that $t : F \circ f_* \to F_\mathcal{K} \circ f_*$ is an isomorphism
for every morphism $f : X' \to X$ whose source is an affine scheme.
\end{lemma}

\begin{proof}
Consider a morphism $f' : X' \to X$ with $X'$ affine. Since the
diagonal of $X$ is affine, we see that $f'$ is an affine morphism
(Morphisms, Lemma \ref{morphisms-lemma-affine-permanence}).
Thus $f'_* : \QCoh(\mathcal{O}_{X'}) \to \QCoh(\mathcal{O}_X)$
is an $R$-linear exact functor
(Cohomology of Schemes, Lemma \ref{coherent-lemma-relative-affine-vanishing})
which commutes with direct sums
(Cohomology of Schemes, Lemma \ref{coherent-lemma-colimit-cohomology}).
Thus $F \circ f'_*$ is an $R$-linear, right exact functor which commutes
with arbitrary direct sums. Whence
$F \circ f'_* = F_{\mathcal{K}'}$ for some $\mathcal{K}'$
on $X' \times_R Y$ by Lemma \ref{lemma-functor-quasi-coherent-from-affine}.
Moreover, given a morphism $f'' : X'' \to X'$ with $X''$ affine
we obtain a canonical identification
$(f'' \times \text{id}_Y)^*\mathcal{K}' = \mathcal{K}''$
by the references already given combined with
Lemma \ref{lemma-functor-quasi-coherent-from-affine-compose}.
These identifications satisfy a cocycle condition given
another morphism $f''' : X''' \to X''$ which we leave it to
the reader to spell out.

\medskip\noindent
Choose an affine open covering $X = \bigcup_{i = 1, \ldots, n} U_i$.
Since the diagonal of $X$ is affine, we see that the intersections
$U_{i_0 \ldots i_p} = U_{i_0} \cap \ldots \cap U_{i_p}$ are affine.
As above the inclusion morphisms
$j_{i_0 \ldots i_p} : U_{i_0 \ldots i_p} \to X$ are affine.
Denote $\mathcal{K}_{i_0 \ldots i_p}$ the quasi-coherent module
on $U_{i_0 \ldots i_p} \times_R Y$ corresponding to
$F \circ j_{i_0 \ldots i_p *}$ as above.
By the above we obtain identifications
$$
\mathcal{K}_{i_0 \ldots i_p} =
\mathcal{K}_{i_0 \ldots \hat i_j \ldots i_p}|_{U_{i_0 \ldots i_p} \times_R Y}
$$
which satisfy the usual compatibilites for glueing. In other words, we obtain
a unique quasi-coherent module $\mathcal{K}$ on $X \times_R Y$
whose restriction to $U_{i_0 \ldots i_p} \times_R Y$ is
$\mathcal{K}_{i_0 \ldots i_p}$ compatible with the displayed identifications.

\medskip\noindent
Next, we construct the transformation $t$. Given a quasi-coherent
$\mathcal{O}_X$-module $\mathcal{F}$ denote $\mathcal{F}_{i_0 \ldots i_p}$
the restriction of $\mathcal{F}$ to $U_{i_0 \ldots i_p}$ and denote
$(\text{pr}_1^*\mathcal{F} \otimes \mathcal{K})_{i_0 \ldots i_p}$
the restriction of $\text{pr}_1^*\mathcal{F} \otimes \mathcal{K}$ to
$U_{i_0 \ldots i_p} \times_R Y$.
Observe that
\begin{align*}
F(j_{i_0 \ldots i_p *}\mathcal{F}_{i_0 \ldots i_p})
& =
\text{pr}_{i_0 \ldots i_p, 2, *}(
\text{pr}_{i_0 \ldots i_p, 1}^*\mathcal{F}_{i_0 \ldots i_p}
\otimes \mathcal{K}_{i_0 \ldots i_p}) \\
& =
\text{pr}_{i_0 \ldots i_p, 2, *}
(\text{pr}_1^*\mathcal{F} \otimes \mathcal{K})_{i_0 \ldots i_p}
\end{align*}
where $\text{pr}_{i_0 \ldots i_p, 2} : U_{i_0 \ldots i_p} \times_R Y \to Y$
is the projection and similarly for the other projection. Moreover, these
identifications are compatible with the displayed identifications
in the previous paragraph. Recall, from Cohomology of Schemes, Lemma
\ref{coherent-lemma-separated-case-relative-cech}
that the relative {\v C}ech complex
$$
\bigoplus 
\text{pr}_{i_0, 2, *}
(\text{pr}_1^*\mathcal{F} \otimes \mathcal{K})_{i_0}
\to
\bigoplus 
\text{pr}_{i_0i_1, 2, *}
(\text{pr}_1^*\mathcal{F} \otimes \mathcal{K})_{i_0i_1}
\to
\bigoplus 
\text{pr}_{i_0i_1i_2, 2, *}
(\text{pr}_1^*\mathcal{F} \otimes \mathcal{K})_{i_0i_1i_2}
\to \ldots
$$
computes $R\text{pr}_{2, *}(\text{pr}_1^*\mathcal{F} \otimes \mathcal{K})$.
Hence the cohomology sheaf in degree $0$ is $F_\mathcal{K}(\mathcal{F})$.
Thus we obtain the desired map
$t : F(\mathcal{F}) \to F_\mathcal{K}(\mathcal{F})$
by contemplating the following commutative diagram
$$
\xymatrix{
&
F(\mathcal{F}) \ar[r] \ar@{..>}[d] &
\bigoplus F(j_{i_0*}\mathcal{F}_{i_0}) \ar[r] \ar[d] &
\bigoplus F(j_{i_0i_1*}\mathcal{F}_{i_0i_1}) \ar[d] \\
0 \ar[r] &
F_\mathcal{K}(\mathcal{F}) \ar[r] &
\bigoplus 
\text{pr}_{i_0, 2, *}
(\text{pr}_1^*\mathcal{F} \otimes \mathcal{K})_{i_0}
\ar[r] &
\bigoplus 
\text{pr}_{i_0i_1, 2, *}
(\text{pr}_1^*\mathcal{F} \otimes \mathcal{K})_{i_0i_1}
}
$$
We obtain the top row by applying $F$ to the (exact) complex
$0 \to \mathcal{F} \to \bigoplus j_{i_0*}\mathcal{F}_{i_0} \to
\bigoplus j_{i_0i_1*}\mathcal{F}_{i_0i_1}$ (but since
$F$ is not exact, the top row is just a complex and not
necessarily exact).
The solid vertical arrows are the identifications above.
This does indeed define the dotted arrow as desired.
The arrow is functorial in $\mathcal{F}$; we omit the details.

\medskip\noindent
We still have to prove the final assertion. Let $f : X' \to X$
be as in the statement of the lemma and let $\mathcal{K}'$
be the quasi-coherent module on $X' \times_R Y$ constructed
in the first paragraph of the proof. If the morphism
$f : X' \to X$ maps into one of the opens $U_i$, then the
result follows from
Lemma \ref{lemma-functor-quasi-coherent-from-affine-compose}
because in this case we know
that $\mathcal{K}_i = \mathcal{K}|_{U_i \times_R Y}$
pulls back to $\mathcal{K}$. In general, we obtain an
affine open covering $X' = \bigcup U'_i$ with $U'_i = f^{-1}(U_i)$
and we obtain isomorphisms
$\mathcal{K}'|_{U'_i} = f_i^*\mathcal{K}_i$ where
$f_i : U'_i \to U_i$ is the induced morphism.
These morphisms satisfy the compatibility conditions needed
to glue to an isomorphism $\mathcal{K}' = f^*\mathcal{K}$
and we conclude. Some details omitted.
\end{proof}

\end{document}
